\chapter{Metodologia de Validação Experimental}\label{cap:metodologia}

Os capítulos anteriores descrevem o que foi construído: a pilha de
software, o simulador e as políticas de controle. Este capítulo descreve o
aparato que decide se um número produzido por esse conjunto pode ser reportado.
São coisas diferentes, e a distinção é o assunto do capítulo.

A prática corrente em aprendizado por reforço aplicado a corrida autônoma
reporta uma métrica agregada --- taxa de conclusão de voltas, tempo de volta,
distância percorrida --- acompanhada de uma descrição do treino
\cite{Betz2022Survey}. O leitor recebe o valor e a receita, e a reprodutibilidade
é tratada como consequência de ambos. O que se argumenta aqui, com base em
falhas medidas e não em princípio geral, é que valor e receita são insuficientes:
entre o experimento e o número existe um conjunto de decisões --- definição do
evento que a métrica conta, protocolo de avaliação, composição da política
avaliada, critério de aceite, momento em que o critério foi fixado --- que não
aparece na receita e que altera o número em magnitude comparável à das diferenças
que o trabalho pretende demonstrar.

Todos os resultados discutidos neste capítulo foram obtidos \textbf{em
simulação}, no ambiente \texttt{f110\_gym} \cite{OKelly2020F1TENTH}. A simulação
é uma escolha de escopo do trabalho, que compara pilotos dentro de um mesmo
simulador e sob um mesmo avaliador, e não uma limitação contornada: o que a
validação descrita garante é que o número comparado entre os pilotos mede o que
o texto diz que mede.

O capítulo está organizado em torno de um argumento único, exposto em sete
partes. A Seção~\ref{sec:val-problema} apresenta quatro falhas reais do projeto
que motivam tudo o que vem depois. A Seção~\ref{sec:val-contrato} descreve o contrato
numérico com identificador, que é a forma escolhida para impedir que uma decisão
de projeto vire folclore. A Seção~\ref{sec:val-congelado} trata de congelar
critérios antes da medição, e de tornar a proveniência do resultado parte do
veredito. A Seção~\ref{sec:val-estatistica} trata da decisão estatística e do
escalonamento automático de amostra. A
Seção~\ref{sec:val-humano} aplica o aparato ao comparativo, com o protocolo
comum, os pilotos humanos e as métricas novas. A Seção~\ref{sec:val-camadas}
descreve as cinco camadas de verificação. A
Seção~\ref{sec:val-limitacoes} declara o que este aparato não cobre.

\section{O problema: um número medido não é um resultado}\label{sec:val-problema}

Quatro episódios do desenvolvimento ilustram o problema melhor do que uma
formulação abstrata. Os quatro produziram números plausíveis, internamente
consistentes e reprodutíveis --- no sentido de que reexecutar o mesmo comando
devolvia o mesmo valor --- e os quatro estavam errados por motivos distintos.

\subsection{O detector de colisão do simulador não detecta
colisão}\label{sec:val-colisao}

O primeiro caso é o mais básico dos quatro, porque não corrompe um número: corrompe
a definição do evento que todos os números contam. Taxa de conclusão de voltas,
tempo de volta e distância percorrida são todos condicionados a ``o veículo não
bateu''. Quem decide isso, no \texttt{f110\_gym}, é uma rotina
(\texttt{check\_ttc\_jit}) que estima tempo até a colisão a partir do próprio scan
do LiDAR. Auditada contra verdade de campo --- o retângulo do veículo amostrado
sobre o mesmo grid de ocupação que o simulador carregou ---, ela apresentou seis
defeitos, todos reproduzidos em simulação e registrados em \texttt{SPEC-COL}.

Três deles bastam para a leitura. O primeiro é geométrico: a rotina dispara dentro
de uma \textbf{faixa} de folga, e não num semiplano; quem atravessa a faixa inteira
num único passo de física não é detectado nem naquele passo nem em nenhum outro
depois, porque a partir dali a folga é negativa e a condição reprova. O segundo é
uma guarda de velocidade (\texttt{if vel != 0.0}) no topo da rotina: veículo parado
encostado no muro não colide. O terceiro é o mais destrutivo, porque atinge o estado
\textbf{físico} e não apenas o que se reporta: a rotina zera as componentes de
estado a partir da quarta, e a quinta é o \textit{yaw}. Toda batida teleporta a
orientação do veículo para 0~rad. Medido: 3,763~rad $\rightarrow$ 0,000.

O que transforma isso de curiosidade em erro de medição é a geometria destes mapas.
As paredes são uma fita de \textbf{9 a 38~cm} com espaço livre atrás, em 40 de 40
seções amostradas nas cinco pistas, enquanto a 8~m/s o veículo percorre 8~cm por
passo de física. Espessura de parede e passo de integração são da mesma ordem: a
travessia num único passo não é o caso patológico, é o caso comum.

O efeito composto foi medido na ponte ROS2, que agravava o quadro por não ler a
chave \texttt{collisions} da observação devolvida pelo simulador --- a marca de
colisão nascia falsa e morria falsa, e não havia tópico que a publicasse. Largando o
veículo a 1,1~m do muro, a 2~m/s, contra ele: o carro atravessou a parede em
$(-44{,}24;\ 3{,}32)$, teve o \textit{heading} zerado no impacto e seguiu dirigindo
sozinho na área externa, parando em
$(-30{,}14;\ 3{,}31)$ --- 14~m adiante do ponto de impacto, fora do circuito. Nada
nessa execução acusou colisão. Uma volta assim é contada como completa, e a
distância percorrida segue sendo integrada do lado de fora da pista.

A correção está em \texttt{SPEC-COL-01} a \texttt{SPEC-COL-06} e redefine o evento
em vez de ajustar o detector: colidiu passa a ser existir ponto do retângulo do
veículo sobre célula ocupada \textbf{ou} a marca do simulador ter acendido; a
checagem é do \textbf{segmento} entre a pose anterior e a atual, e não dos extremos,
que é o que impede atravessar uma parede de 9~cm com as duas poses em espaço livre;
sair dos limites do mapa conta como colisão; o grid consultado é o \textbf{mesmo}
que o simulador usou para gerar o scan, de modo que a geometria da colisão não possa
discordar da geometria que o LiDAR enxerga; e a pose reportada no passo da colisão
preserva o último \textit{heading} válido.

Uma restrição de engenharia acompanha a correção e é parte dela. Ela vive num
invólucro, e o \texttt{f110\_gym} incorporado ao repositório permanece no commit
fixado (\texttt{SPEC-COL-07}), com teste que falha se alguém o alterar. O que o
simulador executa continua sendo o código original, e o scan do ambiente de
pesquisa segue idêntico ao que o mesmo simulador produz no \textit{workspace} ROS2
do repositório --- uma paridade herdada da fase com veículo e ainda verificada
pela Camada~C (Seção~\ref{sec:val-camadas}). Corrigir o critério de colisão sem tocar no simulador é o que
permite as duas coisas ao mesmo tempo.

\subsection{Vazamento de estado entre episódios}

A política da segunda etapa do trabalho (E2, aprendizado por imitação) guarda o
scan do instante anterior e os últimos comandos de esterço e velocidade
(\texttt{scan\_prev}, \texttt{prev\_steer}, \texttt{prev\_speed}) numa closure, e
não expõe método de reinicialização. O avaliador construía \textit{uma} instância
da política e a reutilizava para todos os episódios da avaliação.

A consequência é que o episódio seguinte largava com o scan do instante da
colisão anterior. Do ponto de vista da rede, isso é um deslocamento entre a pose
colidida e a nova largada --- um par de quadros que não existe em lugar nenhum do
treino, onde todo primeiro quadro após o \textit{reset} tem
\texttt{scan\_prev == scan} por construção.

O efeito não é um viés pequeno: é binário. Com duas sementes por pista no
conjunto de generalização, uma pista saía 0/2 ou 6/6, e nunca no meio. E ele se
disfarça de má generalização, que é exatamente a hipótese que um pesquisador
formularia primeiro diante desse padrão. \textbf{Foram oito execuções do E2
reprovadas assim antes de a causa aparecer.} O agravante metodológico é que o
controle não era afetado: o especialista da primeira etapa (E1, \textit{pure
pursuit}) não guarda estado, e por isso era imune --- o que fazia o contraste
entre especialista e rede parecer mérito do especialista.

A correção é uma linha de protocolo, hoje registrada como
\texttt{SPEC-EVAL-01}: cada semente recebe uma política \textbf{recém-construída};
a função de fábrica é chamada uma vez por semente, nunca uma vez por avaliação.
Nenhum estado atravessa a fronteira de episódio.

\subsection{Duas cópias da mesma composição de política}\label{sec:val-duascopias}

A terceira etapa (E3) é um resíduo de aprendizado por reforço somado à saída da
rede do E2. Compor as duas exige montar um vetor de ganhos, aplicar a saída de
uma sobre a da outra e clipar o resultado. Durante um período existiram
\textit{duas} implementações dessa composição: uma usada pelo treino, outra pela
reavaliação.

Elas divergiam num detalhe que qualquer revisor classificaria como inofensivo: o
vetor de ganhos era \texttt{float64} numa e \texttt{float32} na outra. Medido no
\textit{mesmo} ponto de verificação, com as \textit{mesmas} sementes, o resultado
foi \textbf{97,56\%} (40/41 voltas, 1 colisão) contra \textbf{91,43\%} (32/35
voltas, 3 colisões). A diferença induzida na ação é da ordem de $10^{-7}$.

O achado não é que uma das implementações estivesse errada --- nenhuma estava. É
que o sistema amplifica ruído de ponto flutuante ao longo de milhares de passos
de controle. Isso foi medido diretamente, perturbando o vetor de ganhos por
epsilons crescentes e reavaliando o mesmo ponto de verificação com as mesmas
cinco sementes em Spielberg:

\begin{table}[h]
    \centering
    \begin{tabular}{|c|c|c|}
        \hline
        \textbf{$\varepsilon$} & \textbf{E3 (resíduo)} & \textbf{E2 (rede base)} \\
        \hline
        0          & 97,56\%         & 97,56\%         \\
        \hline
        $10^{-12}$ & 93,75\%         & 97,56\%         \\
        \hline
        $10^{-9}$  & 94,29\%         & 97,56\%         \\
        \hline
        $10^{-7}$  & 94,59\%         & 97,56\%         \\
        \hline
        $10^{-5}$  & 97,56\%         & 97,56\%         \\
        \hline
        \textbf{amplitude} & \textbf{0,0381} & \textbf{0,0000} \\
        \hline
    \end{tabular}
    \caption{Completion sob perturbação do vetor de ganhos, mesmo ponto de
    verificação e mesmas sementes (\texttt{SPEC-EVAL-02})}
    \label{tab:val_sensibilidade}
\end{table}

A leitura da Tabela~\ref{tab:val_sensibilidade} tem duas partes. A perturbação
não tem tendência nem ordem: ela reembaralha \textit{qual} episódio bate, não
desloca o resultado numa direção. E o contraste entre as etapas é o achado
central: a política do E2 é \textit{insensível} a perturbações cobrindo sete
ordens de grandeza, enquanto a do E3 espalha 3,8 pontos percentuais. Isso é
coerente com o que as duas fazem --- o resíduo dirige no limite (53,7~s contra
58,8~s da rede base naquela medição), onde uma diferença mínima decide quem bate.

A consequência metodológica está registrada em \texttt{SPEC-EVAL-02} e é
retomada na Seção~\ref{sec:val-limitacoes}: a composição da política avaliada,
\textbf{inclusive tipos numéricos e ordem de operações}, é parte do protocolo, e
não detalhe de implementação. Hoje existe uma única função que monta a
composição, usada pelo treino e pela reavaliação.

\subsection{Três valores para a mesma medição}

A quarta falha não é de código: é de documentação. Três valores diferentes para
a \textit{mesma} medição de distância percorrida numa verificação em tempo real da política ---
337~m, 428~m e 552,7~m --- conviveram simultaneamente no repositório, em
documentos distintos, cada um apresentado como o resultado vigente. Pior: um
deles citava como evidência um arquivo que o contradizia.

Nenhum teste falhava, porque nenhum teste lê prosa. Nenhuma execução divergia,
porque o código estava certo. O defeito vivia inteiramente na camada que o leitor
consulta, e a origem de um dos valores foi posteriormente rastreada: a
verificação não encerrava os nós ROS2 de execuções anteriores, e a leitura de 552~m
saiu com um nó de 1~h~18~min publicando comandos ao lado do nó novo --- duas
políticas comandando o mesmo veículo simulado.

\subsection{O que os quatro casos têm em comum}

Nos quatro, reexecutar o comando devolvia o mesmo número. Nos quatro, o número
estava errado. E nos quatro a distância entre o valor errado e o valor correto é da
mesma ordem de grandeza das diferenças que o trabalho pretende demonstrar --- uma
volta contada como completa depois de o veículo sair do circuito, 6 pontos de
\textit{completion}, uma pista inteira zerada, 200~m de distância percorrida.

Isso delimita o que a validação precisa cobrir. Não basta verificar que o
programa roda, nem que ele é determinístico: é preciso verificar que o evento
contado pela métrica corresponde ao que de fato acontece no mundo simulado, que o
protocolo de avaliação não vaza estado, que existe uma única definição executável
do objeto avaliado, que o critério de aceite foi fixado antes da medição, que o
número citado no texto corresponde a um artefato em disco. Cada uma dessas verificações é uma camada, e é
delas que trata o restante do capítulo. A verificação formal de propriedades de
segurança de controladores neurais \cite{Ivanov2020Verifying} e a busca ativa por
cenários de falha \cite{Bak2021Stress} atacam um problema adjacente e
complementar; o que se discute aqui é anterior a ambos, porque diz respeito à
confiabilidade do próprio número que se pretende verificar.

\section{Contrato numérico com identificador}\label{sec:val-contrato}

O primeiro mecanismo é um documento único, \texttt{SPEC.md}, em que cada decisão
numérica do ambiente de pesquisa é escrita como uma asserção testável com
identificador estável. Na versão vigente, contada em outubro de 2026, são
\textbf{118 asserções}, distribuídas em 20 famílias temáticas, com
identificadores da forma \texttt{SPEC-<FAMÍLIA>-<NN>}.

\subsection{O que uma asserção contém}

Uma asserção não é um comentário: ela fixa um valor, uma unidade e um
comportamento observável. \texttt{SPEC-OBS-01}, por exemplo, estabelece que o
scan de 1080 feixes é reduzido a 60 por média de janelas de 18, e registra a
razão --- $1080/60 = 18$ é exato, a média preserva melhor o obstáculo mais
próximo do que a decimação, e 64 feixes dariam janela fracionária.
\texttt{SPEC-ENV-01} fixa a física em 100~Hz e o controle em 20~Hz, com cinco
passos de simulação por passo de controle. \texttt{SPEC-ACT-02} fixa
$\delta = a_0 \cdot 0{,}4189$~rad, o limite de esterço do modelo
\textit{single-track}.

Além do valor, a maior parte das asserções registra \textbf{a falha que a
originou}. Essa é a diferença entre um documento de especificação e este: quase
todo item nasceu de um defeito medido, e o registro da medição viaja junto com a
regra. \texttt{SPEC-FRE-04} é um exemplo instrutivo. A projeção de Frenet
usada para medir progresso levantava exceção sempre que o deslocamento
longitudinal num passo excedia 2~m, sob a premissa de que todo salto grande fosse
teleporte ou defeito de projeção. A premissa é falsa: Spielberg tem uma forquilha
de raio 0,84~m, menor que a meia-largura de 1,10~m da pista, e no centro de
curvatura todas as normais da linha central se encontram, deixando o pé da
projeção indefinido. Medido: carro a 0,25~m desse centro, \textit{dentro} da
pista, a 4,78~m/s, sem colisão, com salto de 2,39~m em um passo de 50~ms. Nem
levantar nem aceitar em silêncio servem --- levantar mata um treino de doze
minutos por geometria legítima; aceitar dá $+0{,}239$ de recompensa contra
$+0{,}04$ de um passo normal, num ponto \textit{fixo} da pista, que um agente
aprende a explorar. A regra vigente limita o deslocamento a 2~m e \textit{conta}
o evento, com o viés resultante demonstrado conservador. Os três parágrafos que
sustentam essa escolha estão no próprio item do contrato, não numa discussão
perdida em histórico de versão.

\subsection{O teste que impede o contrato de virar prosa morta}

Um documento de especificação sem execução apodrece: alguém acrescenta uma
asserção, ninguém escreve o teste, e o documento passa a descrever um sistema que
não existe. A defesa é a Camada~A, um teste que varre o \texttt{SPEC.md}, extrai
o conjunto de identificadores declarados e o conjunto de identificadores citados
nos arquivos de teste, e \textit{falha} se a diferença for não vazia. A mensagem
de falha é deliberadamente estreita: ou se escreve o teste que falta, ou se
remove o item do contrato.

Há um segundo teste na direção oposta: nenhum teste pode citar um identificador
que não exista no contrato. Sem ele, um erro de digitação numa docstring
satisfaria a cobertura sem verificar coisa alguma.

A ressalva honesta é que essa checagem é \textbf{por ocorrência textual}: ela
garante que o identificador foi citado no arquivo de teste, não que o teste
exerça o que a asserção afirma. A limitação é discutida em detalhe na
Seção~\ref{sec:val-limitacoes}.

\section{Critérios congelados antes da medição}\label{sec:val-congelado}

A segunda família de mecanismos ataca um modo de falha que nenhum teste unitário
alcança: mover a trave depois de ver o número. Ele não exige má-fé --- basta que
o limiar viva num arquivo de configuração editável e que a avaliação possa ser
reexecutada sobre um relatório antigo.

\subsection{Proveniência como veredito, não como metadado}

Duas asserções do contrato não julgam a política: julgam a \textit{medição}. Um
número correto obtido de forma não defensável não é reportável, e o lugar de
dizer isso é o próprio veredito do portão.

\texttt{SPEC-PROV-01} determina que o \texttt{config.yaml} de cada execução
congele, no momento em que a execução nasce, o conjunto de limiares em vigor e
uma assinatura deles (\texttt{gate\_hash}). Se, na hora de avaliar, os limiares
em vigor não baterem com os congelados, o veredito é \texttt{INVALIDO} --- nem
\texttt{PASS} nem \texttt{FAIL}. A motivação está registrada: o contrato afirmava
desde o início que as faixas eram fixadas antes de treinar, e não existia nenhum
artefato capaz de provar isso depois, porque nenhum limiar era gravado na
execução. Era possível editar a configuração e reexecutar a avaliação sobre um
relatório antigo sem que nada denunciasse.

\texttt{SPEC-PROV-02} determina que um \texttt{PASS} obtido com a árvore de
trabalho suja seja rebaixado a \texttt{PASS-NAO-REPRODUZIVEL}. Com mudanças não
commitadas, o identificador de commit registrado não descreve o código que
produziu o resultado --- e esse campo era gravado e nunca lido por ninguém. O
caso que forçou a regra é revelador do quanto o mecanismo precisa ser
automático: uma ação pré-declarada do laço automático disparava a recalibração do
ruído de coleta, e o laudo era gravado num arquivo versionado. A árvore sujava no
meio da execução, e a tentativa seguinte nascia marcada como suja. A quinta
tentativa do E2 passou nos quatro critérios e saiu
\texttt{PASS-NAO-REPRODUZIVEL}: o laço invalidou, pela própria regra, o resultado
que ele mesmo produziu. A correção (\texttt{SPEC-DART-05}) move o laudo novo para
o diretório da execução e para um cache fora do controle de versão; o laudo
citado pelo trabalho passa a ser somente de leitura, atualizado por ato humano
explícito.

\texttt{SPEC-PROV-03} exige que a configuração grave o dispositivo de treino
\textit{resolvido}, nunca o literal \texttt{"auto"}. O dispositivo muda os pesos:
kernels diferentes e ordem de redução diferente produzem resultado diferente com
a mesma semente --- medido, uma divergência de $1{,}34 \times 10^{-2}$ nos pesos
entre CPU e GPU. Sem o campo, uma execução em CPU e outra em GPU sairiam com
configuração byte a byte idêntica e pontos de verificação distintos, e a
configuração é exatamente o artefato que deveria permitir refazer a execução.
Pedir GPU numa máquina sem GPU \textbf{levanta exceção}; cair para CPU em
silêncio devolveria outro experimento sob a mesma configuração.

O valor prático dessas três regras aparece quando o resultado sobrevive a elas. O
tempo do especialista, denominador de todos os critérios de E2 e E3, foi
remedido com a árvore limpa e reproduziu o valor anterior até o último dígito
($T_{pp} = 58{,}99999999994634$~s), inclusive entre arquiteturas distintas
(arm64 e x86\_64). Essa coincidência não era garantida --- aritmética de ponto
flutuante entre ISAs diferentes não precisa concordar bit a bit --- e é o que
autoriza comparar as tabelas novas com as antigas.

\subsection{Ações pré-declaradas por critério}

Congelar o limiar não basta se a escolha do próximo hiperparâmetro for feita
olhando o resultado. O mecanismo correspondente é uma tabela de ações
pré-declaradas, indexada pelo \textit{critério} que reprovou --- não pela etapa
--- e percorrida em ordem fixa pelo laço automático.

A indexação por critério, e não por etapa, foi ela própria uma correção. Na
versão anterior, o laço lia apenas o veredito e descartava a lista de critérios.
Com o critério de generalização reprovando, ele aplicava, em ordem, as cinco
ações declaradas ``do E2'', todas sobre o conjunto de dados de uma única pista, e
nenhuma capaz de mover generalização --- cinco tentativas de aproximadamente uma
hora, garantidamente inúteis. A regra correspondente
(\texttt{SPEC-LOOP-02}) exige que o ajuste venha da lista pré-declarada
\textbf{do critério que reprovou}, na ordem em que os critérios aparecem no
relatório do portão.

Três asserções complementares fecham modos de falha vizinhos.
\texttt{SPEC-LOOP-03} exige que a tabela executável e a tabela textual exibida no
diagnóstico cubram os mesmos critérios com o mesmo número de ações --- quando
eram indexadas de formas diferentes, o diagnóstico prometia uma ordem e o laço
executava outra. \texttt{SPEC-LOOP-07} exige que toda chave declarada seja
parâmetro que a etapa \textbf{de fato consome}, e nasceu de dois casos reais:
\texttt{wall\_w} e \texttt{mu\_std} estavam declarados, existiam nas estruturas
de configuração, e a função construtora do ambiente não os repassava. Um ajuste
que a etapa ignora produz uma execução idêntica à anterior e consome uma
tentativa --- e \texttt{mu\_std} era a \textit{primeira} ação de dois dos quatro
critérios do E3. \texttt{SPEC-LOOP-06} determina que falha das camadas de teste
ou de invariantes não consuma tentativa: é defeito de código ou de ambiente, e se
corrige no código, nunca no hiperparâmetro.

O resultado final do trabalho foi obtido percorrendo essa lista, em ordem, sem
escolher hiperparâmetro olhando o portão. Convém separar quem executou o quê: o
laço automático rodou as quatro primeiras tentativas e parou no
\texttt{PASS-PARCIAL}, que é veredito terminal; os pontos seguintes foram
executados manualmente, seguindo a mesma lista pré-declarada e na ordem dela. A
disciplina é a mesma; o executor não. Os quatro pontos medidos, com uma variável
mudando por vez:

\begin{table}[h]
    \centering
    \begin{tabular}{|p{3.2cm}|c|c|c|p{3.2cm}|}
        \hline
        \textbf{Configuração} & \textbf{Completion} & \textbf{Volta (s)} &
        \textbf{Queda sob $\mu$} & \textbf{Veredito} \\
        \hline
        150k passos                    & 1,0000 & 56,2375 & 0,0000 & PASS-PARCIAL \\
        \hline
        300k passos                    & 0,8793 & 52,4750 & 0,0909 & reprovado (compl., $\mu$) \\
        \hline
        300k + \texttt{wall\_w} 0,2    & 0,9610 & 51,3750 & 0,0667 & reprovado (compl., $\mu$) \\
        \hline
        \textbf{300k + \texttt{wall\_w} 1,0} & \textbf{0,9891} & \textbf{52,5250} &
        \textbf{0,0000} & \textbf{aprovado} \\
        \hline
    \end{tabular}
    \caption{Trajetória do E3 pelas ações pré-declaradas, uma variável por vez}
    \label{tab:val_acoes_e3}
\end{table}

A primeira linha da Tabela~\ref{tab:val_acoes_e3} não é uma reprovação: é a
faixa \texttt{PASS-PARCIAL}, declarada em \texttt{SPEC-GATE-05} \textit{antes}
de treinar para o caso de \textit{completion} aprovada com mediana entre
$0{,}95 \cdot T_{pp}$ e $1{,}00 \cdot T_{pp}$. Registrá-la como ``reprovado''
apagaria uma categoria pré-registrada, que é precisamente o que
\texttt{SPEC-DOC-03} existe para impedir. A tabela é legível como experimento
porque a ordem foi declarada antes. Os dois primeiros pontos desenham uma troca --- dobrar o treino
compra 3,8~s de volta e custa 12 pontos de \textit{completion} --- e os dois
últimos mostram que a troca não era fundamental: era consequência de a recompensa
não cobrar o rasante de muro. A penalidade de colisão só cobra \textit{depois} de
bater, e bater é esparso demais para dar sinal; a penalidade de proximidade cobra
em todo passo em que o rasante acontece. Ligá-la devolveu a \textit{completion} e
ainda melhorou a volta, que é a assinatura de ter mexido no objetivo certo em vez
de andar sobre a curva de compromisso. Se a mesma configuração tivesse sido
encontrada por busca livre, os quatro pontos não sustentariam essa leitura.

\section{Decisão estatística e escalonamento automático}\label{sec:val-estatistica}

Os critérios de aceite na pista de treino são avaliados com 5 sementes
$\times$ 10 voltas e intervalo de confiança de 95\%; as pistas de generalização
usam 5 sementes $\times$ 3 voltas. Três escolhas de método e uma regra de
escalonamento merecem justificativa explícita, porque as quatro afetam diretamente
quais resultados o trabalho pode reportar.

\subsection{Proporção pelo valor pontual, com intervalo ao lado}

A taxa de conclusão de voltas é decidida pelo \textbf{valor pontual}
$\geq$ limiar, com o intervalo de Wilson reportado ao lado
(\texttt{SPEC-GATE-01}). A proporção é calculada sobre o total de ensaios ---
50 voltas em 5 sementes, quando nenhuma semente colide ---, e não sobre as cinco
taxas por semente. O denominador conta as voltas \textit{iniciadas}: como uma
colisão encerra o episódio da semente, uma semente que bate na primeira volta
pesa uma tentativa, e uma que sobrevive às dez pesa dez. O agregado pondera,
portanto, as sementes pelo tempo de sobrevivência, e pode esconder sementes sem
nenhuma volta completa; por isso, onde o agregado não é 1,0, o texto reporta
também a contagem por semente. A escolha de
Wilson em vez do intervalo normal é técnica: ele se comporta bem perto de 0 e de
1, onde a aproximação normal produz limites fora de $[0,1]$, e é exatamente nessa
região que estes portões operam.

A escolha de decidir pelo valor pontual, e não pelo limite inferior do intervalo,
é metodológica e está quantificada no próprio contrato. Exigir limite inferior de
Wilson $\geq$ limiar seria \textbf{inatingível nesta escala}: com 50 voltas e
desempenho perfeito, o limite inferior é 0,929 --- o que reprovaria 95\% e 100\%
sem uma única falha real. Passar em 95\% por esse critério exigiria cerca de 73
voltas perfeitas; passar em 98\%, cerca de 190. Um critério que nenhum resultado
bom consegue atingir não separa bom de ruim: só reprova tudo.

Para a métrica contínua (tempo de volta), o critério é o \textbf{limite superior}
de um \textit{bootstrap} BCa sobre as medianas por semente, com 10.000
reamostragens, $\leq$ limiar (\texttt{SPEC-GATE-02}) --- escolhido porque, com
cinco medianas, o $t$ de Student assumiria uma normalidade que não existe e
produziria intervalo largo demais. Quando o BCa degenera ---
aceleração indefinida quando todos os valores \textit{jackknife} coincidem --- a
implementação cai para o método percentil e \textbf{registra a troca} no campo de
método do relatório, para que o leitor saiba que aquele intervalo é mais fraco.

\subsection{Critérios relativos ao controle}\label{sec:val-relativo}

Um terceiro cuidado aparece no critério de generalização do E2
(\texttt{SPEC-GATE-08}). Um limiar absoluto sobre pistas não vistas seria uma
trave impossível: medido no E1, o próprio especialista faz 0\% em Monza (0/5) e
16,7\% em Oschersleben (1/6). Exigir da rede algo que o controle não faz mede a dificuldade da
pista, não a qualidade da destilação. O critério é, portanto, a razão entre a
taxa da rede e a do controle, por pista, excluindo as pistas onde o controle não
completa nenhuma volta.

Duas salvaguardas foram acrescentadas depois de defeitos medidos. A razão é
\textbf{limitada superiormente em 1,0}, porque sem teto a média mascara: numa
execução anterior à salvaguarda, as razões deram Catalunya 100\%, Oschersleben
300\% e Silverstone 0\%, média 1,333 --- aprovada com uma pista completamente
zerada. Superar o controle numa pista fácil não compra o direito de não andar em
outra. E existe um \textbf{segundo critério, sobre a pior pista}, porque mesmo
com teto a média tolera um zero quando há pistas suficientes.

Há ainda um caso de borda que o portão trata explicitamente: quando o controle
não completa volta em \textit{nenhuma} pista de generalização, o critério devolve
\texttt{INDETERMINADO} e \textbf{não passa}. A versão anterior devolvia aprovação
com limiar zero, o que fazia a ausência de dado virar aprovação automática.

\subsection{O conjunto de generalização se sobrepõe às ações de correção}
\label{sec:val-heldout}

As duas salvaguardas acima tratam de \textit{como} as razões são agregadas.
Nenhuma delas trata de \textit{quais} pistas entram na conta, e é aí que está o
defeito. Ele apareceu ao conferir os artefatos da execução aprovada do E2 contra
o texto que a descrevia, e é da mesma classe da assimetria do escalonamento
discutida na subseção seguinte: um arranjo que \textbf{parece conservador} e não
é. A diferença é que aquela assimetria foi corrigida antes da execução
reportada, e esta só pôde ser corrigida para as execuções seguintes, pelo motivo
exposto ao final da subseção.

O conjunto de pistas de generalização é fixo no código --- Monza, Catalunya,
Silverstone e Oschersleben --- e não depende de onde as demonstrações foram
coletadas. A tabela de ações pré-declaradas da Seção~\ref{sec:val-congelado},
acionada quando o critério de generalização reprova, faz exatamente o oposto:
ela aumenta o conjunto de coleta pista a pista. A primeira ação coleta em
Spielberg e Catalunya; a segunda acrescenta Silverstone; a terceira acrescenta
Oschersleben. \textbf{As duas listas se sobrepunham}, e, quando a execução
aprovada do E2 foi feita, nem o contrato numérico nem o documento de estado
registravam a sobreposição.

A execução aprovada do E2 é a segunda ação: coletou em Spielberg, Catalunya e
Silverstone. Duas das quatro pistas de generalização estavam, portanto, dentro do
conjunto de coleta. A Tabela~\ref{tab:val_heldout} abre por pista o número que o
portão agregou.

\begin{table}[h]
    \centering
    \begin{tabular}{|l|c|c|c|p{3.5cm}|}
        \hline
        \textbf{Pista} & \textbf{Na coleta?} & \textbf{Rede} &
        \textbf{Especialista} & \textbf{Razão pontuada} \\
        \hline
        Catalunya    & \textbf{sim} & 15/15 (100\%)  & 15/15 (100\%) &
        1,00 --- empate \\
        \hline
        Silverstone  & \textbf{sim} & 15/15 (100\%)  & 15/15 (100\%) &
        1,00 --- empate \\
        \hline
        Oschersleben & não          & 15/15 (100\%)  & 1/6 (16,7\%)  &
        1,00 --- teto; a razão medida é $\approx 6\times$ \\
        \hline
        Monza        & não          & 12/13 (92,3\%) & 0/5 (0\%)     &
        excluída --- o controle não completa volta \\
        \hline
    \end{tabular}
    \caption{As quatro pistas do critério de generalização do E2 na execução
    aprovada. O portão pontuou média das razões 1,0000 (limiar 0,80) e pior
    pista 1,0000 (limiar 0,60)}
    \label{tab:val_heldout}
\end{table}

A leitura da Tabela~\ref{tab:val_heldout} é contraintuitiva, e é por isso que ela
está aqui em vez de numa nota de rodapé: \textbf{a sobreposição dilui o critério,
não o infla}. As duas pistas contaminadas contribuem 1,00 cada, mas 1,00 ali é o
empate com o especialista em pista que a rede \textit{viu no treino} --- não mede
generalização nenhuma. As duas pistas genuinamente inéditas são exatamente onde a
rede mais supera o professor, e o critério não dá crédito por nenhuma das duas:
em Oschersleben a razão medida é de aproximadamente seis vezes, e o teto de 1,0 a
corta para 1,00; em Monza a rede completa 12 de 13 voltas contra 0 de 5 do
especialista, e a pista é excluída porque a razão dividiria por zero. O par
1,0000~/~1,0000 que o portão pontuou é, portanto, um \textit{piso} do resultado
medido, não um valor inflado.

A afirmação defensável não é ``generalização 1,00 em três pistas não vistas''. É
esta: a rede completa \textbf{100\% das voltas em Oschersleben e 92\% em Monza},
as duas pistas que nunca entraram na coleta, enquanto o especialista que a
treinou completa 16,7\% e 0\% nessas mesmas pistas. Catalunya e Silverstone, onde
a rede também completa 100\%, entraram na coleta por ação pré-declarada do laço,
e por isso não medem generalização. Uma ressalva de protocolo alcança as duas
pistas inéditas: cerca de 10\% dos pontos da linha de corrida usados como pose
de largada em Oschersleben (10,1\%) e em Monza (10,9\%) já colidem no momento
do \textit{reset} --- em Spielberg, nenhum ---, o que encerra o episódio em um
passo para qualquer política. O efeito deprime por igual a rede e o
especialista nessas pistas, e fica declarado em vez de corrigido porque a
correção mudaria a comparabilidade com todas as execuções anteriores.

Que a coleta é o que move esse número está registrado na sequência das três
execuções, com o conjunto de generalização fixo o tempo todo: com uma única pista
de coleta (Spielberg), a rede não completa volta em nenhuma das quatro; com duas
(mais Catalunya), Oschersleben já sai 12/13 e Silverstone sai 4/8; com três (mais
Silverstone), Oschersleben sai 15/15 e Monza sai 12/13 --- e nem Oschersleben nem
Monza entraram em coleta alguma.

O defeito de desenho, porém, é sobre execuções futuras, e é mais sério que a
diluição. O remédio pré-declarado para reprovar em generalização é
\textbf{coletar nas próprias pistas de generalização}. A terceira ação da lista
absorveria três das quatro --- Catalunya, Silverstone e Oschersleben ---,
deixando Monza como única pista genuinamente inédita; e Monza é justamente a
pista que o critério exclui, porque o especialista não completa volta nela. O
portão continuaria a produzir um número, e esse número continuaria podendo passar
com folga, mas nenhuma das pistas que o compõem seria inédita: o critério ficaria
vazio de conteúdo de generalização \textbf{sem que nada acusasse}. A execução
aprovada parou na segunda ação e não chegou lá --- o que é sorte de ordem de
execução, não propriedade do aparato.

A correção não pode alcançar a execução aprovada, e a razão é metodológica em
vez de técnica: as duas listas eram artefatos congelados antes da medição, e
reescrevê-las para reavaliar aquela execução seria mexer no critério depois de
ver o número, que é exatamente o que \texttt{SPEC-PROV-01} existe para impedir.
O que cabe à execução vigente é o que esta subseção faz --- declarar o que ela de
fato mede. Para as execuções seguintes, a correção foi congelada como regra:
\texttt{SPEC-GATE-09} proíbe que qualquer ação pré-declarada inclua na coleta
uma pista do conjunto de generalização. Como o projeto só dispõe das cinco
pistas já citadas e não há uma sexta pista neutra para coletar, as listas de
ações dos dois critérios de generalização ficaram \textbf{vazias} de propósito:
uma reprovação em generalização passa a encerrar o laço com diagnóstico, em vez
de acionar um remédio que esvaziaria o próprio critério. A execução aprovada
segue, portanto, parcialmente contaminada, e é assim que ela é lida neste
trabalho.

\subsection{Escalonamento por meia largura de intervalo}

A regra que fecha o conjunto é \texttt{SPEC-GATE-07}: se o valor medido cair a
menos de meia largura de intervalo de confiança do limiar, o portão re-executa a
avaliação com 10 sementes antes de declarar \texttt{PASS} ou \texttt{FAIL}. A
justificativa é direta --- nessa faixa, a diferença entre aprovar e reprovar está
dentro do ruído amostral, e decidir ali é decidir por sorteio.

O contrato declara que a regra vale \textit{igualmente} para \texttt{PASS} e
para \texttt{FAIL}, mas não para todas as etapas: só para o E2 e o E3. A
simetria entre \texttt{PASS} e \texttt{FAIL} precisou ser conquistada duas
vezes, e o registro dessa falha é parte do argumento do capítulo. A exclusão do
E1 tem motivo próprio, discutido ao final desta seção.

A primeira implementação chamava o escalonamento apenas na segunda etapa. A
segunda o condicionava a \texttt{ok and iv.low < threshold}, isto é, só remediava
quem já havia passado. O efeito é que o escalonamento sabia transformar
\texttt{PASS} em \texttt{FAIL} e nunca o contrário. Nenhuma das duas versões
produz um resultado inflado --- e é exatamente por isso que passam despercebidas:
uma regra assimétrica nessa direção \textbf{parece conservadora}. Na prática, ela
é um limiar mais duro do que o declarado, aplicado sem que o texto do critério o
diga, o que é uma forma silenciosa de mover a trave.

O que expôs a assimetria foi um padrão: quatro tentativas seguidas do E3
reprovadas em 97,6--97,8\% contra um limiar de 98\%, sempre por \textit{uma}
colisão em aproximadamente 45 voltas, com a distância ao limiar (0,0022)
\textbf{vinte vezes menor} que a meia largura do intervalo (0,0559). O critério
de tempo de volta já usava a versão simétrica; a assimetria estava apenas no de
\textit{completion}.

Corrigida a regra, ela passou a disparar sozinha em ambas as direções. Um ponto
de verificação do E2 que media 100\% com 5 sementes foi re-medido com 10 e
encontrou uma colisão: 0,9896 é o número honesto, e foi a primeira vez que a
regra rodou numa execução real. Os dois resultados finais mostram a regra nas
duas funções que ela tem. O E3 aprovado \textit{foi} escalonado, e o relatório
do portão registra a re-execução com 10 sementes. O E2 aprovado \textit{não} foi
escalonado, e não deveria ter sido: com valor 1,0000 contra
limiar 0,95, a distância ao limiar (0,05) excede a meia largura do intervalo
(0,0357), e a condição não disparou. Escalonamento é remédio num caso e filtro no
outro.

Resta o E1, que o contrato exclui da regra. No código, apenas o E2 e o E3 passam
pelo invólucro que re-executa com mais sementes; o E1 chama a avaliação do
portão diretamente. O motivo é que o limiar do E1 é 1,00, o teto da própria
escala: um resultado perfeito tem distância zero até ele, e a condição de meia
largura avalia verdadeiro em \textit{toda} execução, sem separar nada. Ligá-la
teria ainda um custo medido: com 100 voltas, uma única colisão dá 0,9900 e
reprova o E1, que fixa o $T_{pp}$ de que dependem todos os critérios do E2 e do
E3 --- remedir automaticamente uma medição limítrofe por construção, e não por
sorte de amostragem, arriscaria trocar uma linha de base boa por outra pior sem
sinal de que algo mudou. O efeito é visível no artefato: o relatório do E1
aprovado traz o critério de \textit{completion} marcado como limítrofe, com 5
sementes e \textbf{nenhuma nota de re-execução}, que é o comportamento
declarado. Convém registrar a cronologia: o código sempre se comportou assim,
mas a justificativa foi escrita no contrato em setembro de 2026, depois da
medição do E1 --- é uma exclusão declarada, não um pré-registro. A consequência
para a leitura dos resultados está na Seção~\ref{sec:val-limitacoes}.

\subsection{Uma taxa estimada com poucos eventos, e o que torna a escolha
robusta}\label{sec:val-poisson}

A escolha de método que falta não está num critério de aceite: está na calibração do
ruído de coleta do E2, e é discutida por último porque é o ponto do trabalho em que
uma decisão se apoia na estimativa estatisticamente mais fraca de todas. O
procedimento, descrito no Capítulo~\ref{cap:estArte}, percorre uma grade de
níveis de ruído em ordem decrescente e escolhe, por pista, o maior nível cujo risco
medido cabe num alvo de colisões por quilômetro fixado antes de qualquer medição
(\texttt{SPEC-DART-02}). Se nenhum nível couber, usa-se o menor da grade e o laudo
registra a escolha como não atingida --- calibração que não alcançou o alvo não é
calibração silenciosamente bem-sucedida.

O ponto fraco é o tamanho da amostra. A medição de um nível se encerra ao
atingir cinco colisões, o que nos níveis mais ruidosos acontece em menos de
5~km, e os níveis de risco baixo rodam entre cerca de 12,5 e 21~km de condução.
Mesmo esses produzem \textbf{pouquíssimos eventos}: no nível escolhido para
Spielberg, exatamente uma colisão em 13,5~km. Contagens dessa ordem carregam
incerteza de Poisson larga --- um único evento sustenta um intervalo de 95\% para a
taxa que vai de cerca de 0,002 a cerca de 0,4 colisões por quilômetro, duas ordens
de grandeza. Tomada isoladamente, a taxa estimada em cada nível não é um número
publicável.

O que torna a \textbf{escolha} robusta apesar disso é a granularidade da grade
perto do joelho da curva, e isso está medido. Em Spielberg, o nível 0,02~rad dá
0,074 colisões/km e o nível vizinho 0,03~rad dá 1,091 --- um fator de quinze entre
níveis adjacentes, muito maior que a largura do intervalo que separa cada um deles
de zero. A decisão que o procedimento toma não é ``qual é o risco'', e sim ``este
nível cabe no alvo ou não cabe''; a distância entre níveis é grande o bastante para
que o ruído amostral não altere essa classificação. O alvo classifica níveis; ele
não é uma estimativa publicável do risco absoluto, e nenhum resultado do trabalho o
usa como tal.

Declarar isso importa por simetria com o resto do capítulo: o mesmo trabalho que
recusa decidir um portão de \textit{completion} dentro do ruído amostral --- o
escalonamento por meia largura de intervalo da subseção anterior --- decide a
calibração de ruído com uma colisão por medição. As duas condutas são
compatíveis porque as perguntas são diferentes --- uma estima um valor contra um
limiar próximo, a outra ordena níveis separados por um fator de quinze ---, mas a
diferença precisa estar escrita, e não deduzida pelo leitor.

\section{Protocolo comum, pilotos humanos e métricas do comparativo}\label{sec:val-humano}

As seções anteriores fixam como um número é medido e aceito. Esta seção aplica o
mesmo critério ao comparativo entre pilotos: o protocolo é escrito e congelado
antes da coleta, e vale para todos os pilotos, automáticos e humanos.

\subsection{Protocolo comum}

Todos os pilotos correm em Spielberg, com randomização do atrito e do ruído do
LiDAR e largada em ponto aleatório da linha de corrida, sob o mesmo avaliador e os
mesmos limites de ação (esterço de $\pm 0{,}4189$~rad, velocidade de 1 a 8~m/s,
atraso de esterço de 100~ms e passo de 0,05~s). Uma colisão encerra o episódio, e
voltas não iniciadas não contam como tentadas. Os pilotos automáticos correm cinco
\textit{seeds} (1000 a 1004) de dez voltas. Os humanos correm três sessões
de cinco voltas, nas \textit{seeds} 1000, 1001 e 1002, que são as três primeiras
dos automáticos: largada e atrito coincidem, e a comparação é pareada.

\subsection{Protocolo dos pilotos humanos}

Em cada vista, o piloto faz 10~minutos de treino, medidos no relógio, e em seguida
três sessões que contam. O treino usa \textit{seeds} próprias (9000, 9001, \dots),
reinicia o carro após uma colisão, é gravado e não entra nos resultados, exceto na
curva de aprendizado. Durante o treino, o pesquisador que conduz a sessão responde apenas a dúvidas sobre
os controles, sem sugerir trajetória, velocidade ou estratégia. Cada sessão tem
contagem regressiva de 3~s e termina ao completar cinco voltas ou ao primeiro
choque; as voltas completas até ali contam, e a volta em andamento não. Nenhuma
sessão é repetida por desempenho ruim. Uma sessão só é excluída por motivo externo
ao desempenho (falha de software, falha do \textit{joystick}, interferência na
sala ou pedido do voluntário), com o motivo registrado, e a decisão não olha o
resultado da volta.

A ordem das vistas é fixa: primeiro a pista inteira, depois só o LiDAR, na mesma
visita e com pausa de pelo menos 5~minutos entre elas. As sessões são presenciais,
no mesmo computador, monitor e \textit{joystick}. Participam os dois autores e de
três a cinco voluntários, e os dois grupos são reportados em separado.

\subsection{Critérios congelados}

O protocolo recebe um \textit{hash}, calculado pela mesma técnica do \textit{hash}
dos critérios de aceite e gravado no registro de cada piloto e vista. Estão
congelados a pista, a randomização e a largada; as \textit{seeds} de sessão e de
treino; o número de sessões, de voltas e a duração do treino; a ordem e o conteúdo
das vistas; o mapeamento do \textit{joystick}, a zona morta e os limites de ação;
o alcance de visualização do LiDAR e a contagem regressiva. Mudar qualquer um
exige refazer a coleta. Antes do congelamento, uma sessão-piloto com um dos autores
valida o roteiro e não entra nos resultados.

\subsection{Ética}

O voluntário lê e assina o Termo de Consentimento Livre e Esclarecido (TCLE) em
duas vias e preenche uma ficha curta com a experiência em jogos de corrida
(nenhuma, alguma ou bastante). A ficha não tem nome, e nos dados, nos nomes de
pasta e nas mensagens de versionamento aparece apenas um código anônimo (P01,
P02, \dots). O TCLE, que tem o nome, fica em papel e separado dos dados. O termo
segue as Resoluções CNS 466/2012 e 510/2016. Se o projeto for submetido ao Comitê
de Ética em Pesquisa, nenhuma coleta com voluntários ocorre antes da aprovação;
se o parecer não sair a tempo, o estudo se limita aos autores e a participação de
voluntários passa a trabalho futuro. O voluntário pode interromper a qualquer
momento, e a retirada do consentimento implica descartar os seus dados.

\subsection{Métricas}

As métricas por piloto são calculadas sobre os dados gravados passo a passo.
\begin{itemize}
    \item \textbf{Conclusão de volta}, como proporção de voltas completas entre as
          tentadas, com intervalo de Wilson, como nas seções anteriores.
    \item \textbf{Tempo de volta}, pela mediana e pelo melhor valor entre as voltas
          completas.
    \item \textbf{Colisões}, em contagem.
    \item \textbf{Taxa de esterço}, a média e o percentil 95 de
          $|\Delta\delta|/\Delta t$ em rad/s, com $\delta$ o esterço comandado e
          $\Delta t = 0{,}05$~s, como medida de suavidade do comando.
    \item \textbf{Desvio da linha de corrida}, a média e o percentil 95 da
          distância da pose à linha de corrida, densificada para espaçamento de
          no máximo 0,05~m antes da medição, para que o erro de discretização
          não entre como desvio do carro.
    \item \textbf{Perfil de velocidade}, a velocidade média por faixa de 1~m ao
          longo da volta.
    \item \textbf{Curva de aprendizado}, o tempo de cada volta em ordem cronológica,
          nas fases de treino e de cada uma das três sessões.
\end{itemize}
A volta é segmentada pelo progresso acumulado, e se foi completa e qual foi o seu
tempo vêm sempre do resultado da volta no avaliador, nunca de inferência sobre os
passos gravados. Os agregados usam apenas voltas completas.

\subsection{Estudo de caso descritivo}

Com poucas pessoas, sem amostragem aleatória, o estudo é descritivo. Reporta-se,
por piloto e vista, a tabela das métricas acima, a faixa do grupo humano (melhor,
mediana e pior) e a posição de cada política em relação a ela. Não há teste de
hipótese entre humanos e máquinas, porque a amostra não o sustentaria e porque
ele sugeriria uma generalização que o desenho não permite.

\subsection{Reavaliação dos pilotos aprendidos e do especialista}

O E1, o E2 e o E3 são reavaliados a partir dos \textit{checkpoints} aprovados, sem
retreino, pelo protocolo comum, com randomização e com o gravador passo a passo
(cinco \textit{seeds} de dez voltas). O E1 oficial foi medido \textbf{sem}
randomização (59,000~s, 50 de 50 voltas), e o E1 do comparativo é medido com ela:
são dois números distintos, que convivem no texto com a indicação de qual é qual.
Os números oficiais de E1, E2 e E3 continuam sendo os dos critérios congelados
anteriores; a reavaliação não os substitui.

\section{As cinco camadas de validação}\label{sec:val-camadas}

Os mecanismos das seções anteriores estavam originalmente espalhados: a suíte de
testes de um lado, um subcomando de invariantes de outro, e o portão científico
só dentro de uma etapa. Quem quisesse saber se o repositório estava saudável
precisava lembrar de três comandos e interpretar três saídas. As cinco camadas
consolidam isso num comando e num veredito, e a Tabela~\ref{tab:val_camadas}
resume o que cada uma cobre. As letras nomeiam exclusivamente as camadas do
verificador (\texttt{SPEC-DOC-08}); o laço automático e o portão científico não
são camadas e não recebem letra.

\begin{table}[h]
    \centering
    \begin{tabularx}{\textwidth}{|c|l|X|}
        \hline
        \textbf{} & \textbf{Camada} & \textbf{O que verifica} \\
        \hline
        A & specs & Todo identificador do contrato é citado por algum teste \\
        \hline
        B & testes & Três suítes: ambiente de pesquisa, raiz do repositório
        (servidor MCP e coletor de dados) e um terceiro workspace do repositório, este sem os
        testes dependentes de dispositivos externos \\
        \hline
        C & invariantes & Paridade com o \textit{workspace} ROS2 do repositório,
        \textit{hash} do mapa, commit do simulador vendorizado \\
        \hline
        D & documentação & Nenhum documento declara arquivo, identificador ou
        contagem que não existe \\
        \hline
        E & proveniência & Limiares congelados antes do treino, árvore de
        trabalho limpa \\
        \hline
    \end{tabularx}
    \caption{As cinco camadas de validação e seu escopo}
    \label{tab:val_camadas}
\end{table}

Um princípio de projeto atravessa as cinco: \textbf{uma camada que explode conta
como falha, não como interrupção}. Quando se está diagnosticando um repositório
quebrado, o valor está justamente em ver todas as camadas de uma vez, e um
\textit{check} que levanta exceção levaria junto os outros quatro. O mesmo vale
dentro da Camada~C, onde um invariante que não consegue ser avaliado reporta
falha \textbf{com o motivo}, em vez de propagar a exceção
(\texttt{SPEC-CHK-02}) --- o modo de falha que a camada existe para pegar é
justamente o silencioso.

\subsection{Camadas A a C}

A Camada~A é a cobertura do contrato, descrita na Seção~\ref{sec:val-contrato}. A
Camada~B roda três suítes de teste: a do ambiente de pesquisa; a da raiz do
repositório, que cobre o servidor de ferramentas e o coletor de dados e tem
outro conjunto de dependências; e a de um terceiro workspace do repositório, executada com a exclusão dos
testes marcados como dependentes de dispositivos externos, de desempenho ou lentos. Na
contagem de outubro de 2026, são 433 testes aprovados na primeira, 160 na
segunda e 551 na terceira, com 2 pulados e 24 excluídos pelos marcadores.
Mantê-las sob o mesmo comando é o que impede uma delas de voltar a ficar
invisível --- uma suíte anterior existiu por meses sem que nada a executasse,
até desaparecer. Testes aprovados verificam lógica no computador de
desenvolvimento, em simulação.

A Camada~C verifica invariantes de integração: paridade numérica entre a
inferência no ambiente de pesquisa e no \textit{workspace} ROS2 do repositório,
integridade do mapa por \textit{hash}, e alinhamento do commit do simulador
vendorizado. O \textit{workspace} ROS2, executado num ambiente isolado próprio, é
herança da fase com veículo (Capítulo~\ref{cap:proposta}); nenhum número deste
texto é medido nele, e o invariante de paridade fica apenas como garantia de que o
mapa e a inferência do ambiente de pesquisa coincidem com os dele. São seis invariantes, com a
numeração pulando dois identificadores descartados; o contrato registra
explicitamente que os identificadores \textbf{não são reaproveitados}, porque
renumerar tornaria ilegíveis os relatórios já gravados
(\texttt{SPEC-CHK-01}). A camada roda \textbf{antes de cada tentativa} do laço
automático (\texttt{SPEC-CHK-03}): sem isso, uma aprovação podia sair com o mapa
divergente ou o commit do simulador desalinhado --- números
plausíveis medidos na pista errada.

\subsection{Camada D: documentação verificável}

A Camada~D é a resposta ao terceiro caso da Seção~\ref{sec:val-problema}. Ela
varre recursivamente todos os documentos do repositório e verifica quatro
classes de afirmação: caminho de arquivo citado entre crases, identificador do
contrato, contagem de testes declarada ao lado do nome de um arquivo de teste, e
veredito de etapa.

A quarta classe merece destaque, porque é ela que delimita o que a camada
consegue e o que não consegue. \texttt{SPEC-DOC-03} exige que todo veredito de
etapa enunciado num documento (\texttt{PASS}, \texttt{FAIL},
\texttt{PASS-PARCIAL}, \texttt{PASS-NAO-REPRODUZIVEL}, \texttt{INVALIDO}) cite,
na mesma seção, o identificador da execução que o produziu, e coincida com o
veredito gravado no relatório de portão daquela execução --- o contrato é a única
exceção, porque define as faixas em vez de afirmar resultado. O contraste com
\texttt{SPEC-DOC-02} é o argumento inteiro da camada: julgar número medido em
prosa exigiria um oráculo do que é ``a medição atual'', que não existe; para
\textit{veredito}, o oráculo existe e está em disco. O caso que motivou a regra é
o mesmo modo de falha da Seção~\ref{sec:val-problema}, aplicado a uma categoria
em vez de a um número: um documento afirmou \texttt{PASS-PARCIAL} para o E3
durante duas semanas depois de a execução vigente ter saído \texttt{PASS}, e
nenhuma camada discordava.

O escopo é estreito de propósito, e o contrato registra por quê: um verificador de
prosa daria falso positivo e seria desligado na primeira semana. Três decisões de
projeto vieram de ruído medido. A varredura passou a ser recursiva sobre o
repositório inteiro porque o glob de topo era um buraco silencioso --- documentos
em subdiretórios citavam scripts inexistentes e a camada devolvia ``nenhuma
afirmação falsa'' porque nunca abria o arquivo (\texttt{SPEC-DOC-01}). Caminhos
contendo marcadores de padrão (\texttt{<...>}, \texttt{*}, \texttt{AAAA-MM-DD})
deixaram de ser cobrados, porque descrevem uma forma e não um lugar; eram 87 dos
302 problemas reportados antes da distinção (\texttt{SPEC-DOC-04}). E os
problemas passaram a ser reportados pelo caminho relativo à raiz, não pelo nome
do arquivo, porque o repositório tem mais de trinta arquivos chamados
\texttt{README.md} e a mensagem antiga não permitia saber qual documento mentia
(\texttt{SPEC-DOC-06}).

Dois diretórios são explicitamente ignorados, com justificativa registrada: o
scratch de processo das ferramentas de automação, que é rascunho e não
documentação do projeto (\texttt{SPEC-DOC-05}), e uma cópia vendorizada de um hub
externo de agentes, cujos caminhos citados descrevem a estrutura do repositório
de origem e são referências condicionais por desenho --- aproximadamente 128 dos
problemas reportados (\texttt{SPEC-DOC-07}).

\subsection{Camada E: proveniência}

A Camada~E aplica, sobre o repositório como um todo, as regras descritas na
Seção~\ref{sec:val-congelado}. Ela verifica que a árvore de trabalho está limpa e
que a execução vigente de cada etapa --- a mais recente aprovada --- congelou
limiares que ainda batem com os de hoje. Execuções antigas aprovadas sob limiares
diferentes são fato histórico, não problema: o que precisa ser defensável é o
resultado que o repositório reporta \textit{hoje}.

\section{Limitações declaradas}\label{sec:val-limitacoes}

Um aparato de validação que não declara o que não cobre pede ao leitor a mesma
confiança cega que ele pretende substituir. Esta seção lista as limitações
conhecidas, com o custo de cada uma quantificado sempre que houve medição.

\subsection{A Camada A verifica citação, não exercício}

A cobertura do contrato é verificada por ocorrência textual do identificador nos
arquivos de teste. Isso garante que o identificador foi \textit{citado}, não que
o teste \textit{exerça} o que a asserção afirma. Um teste que cite
\texttt{SPEC-ENV-01} numa docstring e verifique outra coisa satisfaz a camada.

A limitação é estrutural: decidir automaticamente se um teste exerce uma asserção
escrita em português exigiria um oráculo que não existe. O que a camada garante
--- que nenhuma asserção do contrato ficou órfã e que nenhum teste cita
identificador inventado --- é menos do que se gostaria, e é mais do que a
ausência da camada garantiria. A defesa remanescente é revisão humana do par
asserção/teste, e ela não é automatizada.

\subsection{A Camada D não julga número medido em prosa}

A Camada~D verifica caminho citado, identificador do contrato e contagem de
testes. Ela \textit{não} julga número medido escrito em prosa, e isso é limitação
conhecida, não esquecimento: julgar prosa exigiria um oráculo do que é ``a
medição atual'', que não existe.

O custo dessa limitação está documentado e é o caso da
Seção~\ref{sec:val-problema}: 337~m, 428~m e 552,7~m para a mesma medição de
distância percorrida conviveram no repositório, um dos documentos citando como evidência um
arquivo que o contradizia. Nenhuma camada acusou.

Um exemplo residual sobrevive no próprio registro do projeto e fica declarado
aqui: a afirmação ``nos seis conjuntos de dados, 389 milhões de leituras'' não
fecha consigo mesma --- 389 milhões é exatamente \textit{um} conjunto de dados
($360.417 \times 1080 = 389.250.360$), não seis. A afirmação está marcada como
inconsistente no documento de estado, por ato humano, porque nenhuma camada é
capaz de marcá-la sozinha.

Um segundo caso, do mesmo tipo, apareceu durante a redação deste capítulo e já
foi corrigido: a prosa de \texttt{SPEC-GATE-08} registrava que o especialista
fazia 33\% em Oschersleben, enquanto o relatório de avaliação do E1 em disco
registra 1/6, isto é 16,7\%. O contrato foi corrigido para o valor do artefato,
com a nota de que o número anterior não tinha artefato que o sustentasse, e a
Seção~\ref{sec:val-estatistica} usa o mesmo valor. Nenhuma camada acusou a
divergência enquanto ela existiu, pela razão exata desta subseção --- é número
medido escrito em prosa --- e com um agravante: ela estava no próprio contrato,
que é o único documento isento da checagem de veredito, justamente por definir
faixas em vez de afirmar resultado. A isenção é correta para veredito e não
protege número em prosa; a correção veio de leitura humana.

\subsection{O limiar cai dentro do intervalo nos três portões aprovados}

Esta é a ressalva estatística mais importante do trabalho, e ela vale para os
três resultados aprovados. Como o critério de \textit{completion} é o valor
pontual (Seção~\ref{sec:val-estatistica}), o intervalo de Wilson reportado ao
lado \textbf{contém o limiar} nos três casos:

\begin{table}[h]
    \centering
    \begin{tabular}{|c|c|c|c|c|}
        \hline
        \textbf{Etapa} & \textbf{Valor} & \textbf{Limiar} & \textbf{IC95 de Wilson} &
        \textbf{Contém o limiar?} \\
        \hline
        E1 & 1,0000 & 1,00 & [0,9287; 1,0000] & sim \\
        \hline
        E2 & 1,0000 & 0,95 & [0,9287; 1,0000] & sim \\
        \hline
        E3 & 0,9891 & 0,98 & [0,9410; 0,9981] & sim \\
        \hline
    \end{tabular}
    \caption{O limiar cai dentro do intervalo de confiança nos três portões
    aprovados}
    \label{tab:val_ic}
\end{table}

Dito de outro modo: na Tabela~\ref{tab:val_ic}, em nenhum dos três casos os dados
\textit{excluem} a hipótese de
que a taxa verdadeira esteja abaixo do limiar. Isso precisa ser dito, e não é um
defeito descoberto depois --- é consequência de uma decisão pré-registrada, cujo
motivo está quantificado no contrato e foi escrito antes de qualquer medição:
exigir o limite inferior de Wilson tornaria o critério inatingível nesta escala,
já que 50 voltas perfeitas dão limite inferior 0,929, e passar em 95\% exigiria
cerca de 73 voltas perfeitas, em 98\%, cerca de 190. O orçamento de simulação do
trabalho não comporta esse número de voltas, e um critério que nenhum resultado
bom atinge não separa bom de ruim.

Três observações limitam o alcance da ressalva; uma quarta a torna mais aguda
numa das linhas. O resultado
aprovado do E3 já é um valor escalonado para 10 sementes pela regra de meia
largura de intervalo, e não a primeira medição. A mais severa das reprovações do
E3 não depende desta questão: 0,8793 com 10 sementes tem o intervalo inteiramente
abaixo do limiar --- ao contrário da reprovação seguinte, 0,9610, cujo intervalo
o contém tanto quanto o da aprovação, e que a Seção~\ref{sec:res-e3} abre por
esse motivo. E os critérios de robustez a $\mu$
do E3 deram 0,0000 e 0,0019 contra limiares de 0,03 e 0,05, onde não há
ambiguidade alguma.

Convém notar que ``o limiar cai dentro do intervalo'' e ``a regra de
escalonamento dispara'' \textbf{não são o mesmo teste}, e no E2 eles discordam. O
intervalo de Wilson de 50 acertos em 50 ensaios é assimétrico em torno do valor
pontual: 0,95 cai dentro de $[0{,}9287;\ 1{,}0000]$, enquanto a distância ao
limiar (0,05) excede a meia largura do intervalo (0,0357) --- e o portão,
corretamente, não escalonou. A ressalva desta tabela é, portanto, mais
conservadora que o critério operacional do portão, e é assim que ela deve ser
lida.

A ressalva é mais aguda na linha do E1, que é a linha para a qual o leitor olha
primeiro, porque valor e limiar coincidem em 1,00. O critério foi marcado como
limítrofe pelo próprio portão --- a condição de \texttt{SPEC-GATE-07} avaliou
verdadeiro --- e \textbf{não foi remedido}, porque, como descrito na
Seção~\ref{sec:val-estatistica}, o contrato exclui essa etapa do escalonamento:
com o limiar no teto da escala, a condição é verdadeira em toda execução e
remedir só arriscaria a linha de base. O E1 permanece medido em 5 sementes e 50
voltas, com o limiar dentro do intervalo. Isso não move o $T_{pp}$, que é uma
mediana de tempo de volta e não uma proporção, e que reproduziu dígito a dígito
entre duas arquiteturas; mas o leitor deve saber que a ausência de re-execução é
decisão declarada, justificada no contrato depois da medição, e não uma
propriedade estatística do resultado.

A conduta adotada no texto é a única defensável: reportar o valor pontual
\textbf{com o intervalo ao lado}, declarar que o critério é o valor pontual por
decisão pré-registrada, e declarar o porquê. Esconder o intervalo seria erro;
trocar o critério depois de ver o número seria pior, e é exatamente o que
\texttt{SPEC-PROV-01} existe para impedir.

\subsection{O intervalo de Wilson não vê a segunda fonte de variabilidade}

O intervalo de Wilson supõe ensaios de Bernoulli independentes e mede a
variabilidade \textbf{entre episódios}. A Tabela~\ref{tab:val_sensibilidade}
exibe uma segunda fonte, de natureza diferente: a realização numérica da política
composta. Para o E3, a amplitude sob perturbação irrelevante (0,0381) é da mesma
ordem que a distância até o limiar, e o intervalo de confiança reportado não a
enxerga.

A consequência para a leitura dos resultados é explícita: onde a amplitude sob
perturbação irrelevante for comparável à distância até o limiar, o número
reproduzível é a \textit{distribuição}, não a realização. Quem reimplementar a
composição --- outra ordem de operações, outro tipo numérico, outra versão da
biblioteca de álgebra --- vai medir outro valor da faixa. Para o E2, com
amplitude nula sob as mesmas perturbações, o resultado não sofre a mesma
ressalva.

\subsection{Parte da evidência histórica não é mais auditável}

Os diretórios de execução não são versionados --- ocupam 4,2~GB --- e foram
reconstruídos ao migrar o projeto de uma máquina arm64 para o host x86\_64 atual.
As execuções da máquina anterior não sobreviveram. Os números medidos nelas
continuam registrados porque são história real do projeto, mas \textbf{não têm
artefato para reconferir}, e isso está declarado item a item no documento de
estado: a tabela dos seis pontos de verificação reavaliados, o histórico de
reavaliações e o conjunto de dados de coleta correspondente estão ausentes.

Isso alcança a Tabela~\ref{tab:val_sensibilidade} deste capítulo. As
execuções que sustentam o estudo de sensibilidade a perturbação, e o par
97,56\%~/~91,43\% que abre a Seção~\ref{sec:val-problema}, são da máquina
anterior e estão entre as ausentes. O argumento que elas sustentam --- que o
resíduo amplifica ruído numérico e a rede base não --- continua sendo o melhor
disponível, mas o leitor precisa saber que ele repousa sobre medição registrada,
não sobre artefato reconferível, e que reproduzi-lo exigiria refazer o estudo
nesta máquina.

O que \textit{é} auditável são as medições de portão e de avaliação feitas no
host atual: as Tabelas~\ref{tab:val_acoes_e3} e \ref{tab:val_ic} batem dígito a
dígito com os relatórios em disco, e o tempo do especialista da máquina anterior
é confirmável pelo histórico de versão de um arquivo de ativo. As
medições de execução em tempo real feitas em fases anteriores do projeto não
têm essa propriedade: viviam em relatórios de verificação sobrescritos a cada
execução, que não existem mais, e por isso não sustentam nenhum número deste
texto.

\subsection{O escopo é a simulação, por escolha}\label{sec:val-escopo-simulacao}

O trabalho compara pilotos --- clássicos, de aprendizado de máquina e humanos
--- dentro de um mesmo simulador e sob um mesmo avaliador. Por isso todos os
resultados deste capítulo e do seguinte foram obtidos em simulação, e essa é uma
escolha de escopo: comparar sob condições idênticas e repetíveis exige controlar
atrito, largada e ruído de sensor, e é o simulador que oferece esse controle.
Disso decorre o que as medições não cobrem: efeitos de atrito real, de latência
de atuadores e de ruído de sensor real não são avaliados, e nenhuma conclusão
deste texto os atinge.

O aparato descrito neste capítulo, por sua vez, não depende do simulador
escolhido: as camadas, o contrato numérico e a proveniência valem para qualquer
ambiente em que o critério de aceite possa ser congelado antes da medição.

\subsection{Limitações do comparativo}\label{sec:val-limitacoes-comparativo}

\begin{itemize}
    \item \textbf{Ordem fixa das vistas.} Todos fazem primeiro a pista inteira e
          depois só o LiDAR. O efeito de aprendizado e o de cansaço ficam misturados
          com o efeito da vista, e tendem a favorecer a segunda. Não houve
          contrabalanço, pela amostra pequena; a curva de aprendizado mostra o
          tamanho do efeito de prática.
    \item \textbf{Informação desigual na vista da pista inteira.} O humano vê o
          mapa completo e pode antecipar curvas, enquanto a política tem apenas o
          LiDAR do passo atual. A vista só LiDAR é a comparação em informação
          equivalente.
    \item \textbf{Amostra pequena.} Dois autores e de três a cinco voluntários, sem
          amostragem aleatória: os resultados descrevem esse grupo e não permitem
          generalizar para pilotos humanos.
    \item \textbf{Autores separados dos voluntários.} Os autores conhecem o
          simulador, a pista e o estudo, e um deles tem bastante experiência em jogos
          de corrida. Por isso o grupo é um campo dos dados e os resultados são
          mostrados por grupo.
    \item \textbf{Uma só pista, um só \textit{joystick}, um só computador.} Spielberg
          é a única pista do comparativo.
    \item \textbf{Colisão encerra a sessão.} O número de voltas humanas fica
          reduzido para quem bate cedo, o que alarga o intervalo de Wilson e
          limita a comparação do tempo de volta a voltas completas.
    \item \textbf{Treino curto e atraso do programa.} Dez minutos por vista podem não
          estabilizar o desempenho, e o passo é segurado até 0,05~s de relógio, com
          os passos que ultrapassam o prazo medidos e relatados.
    \item \textbf{Dois números de E1.} O E1 oficial, sem randomização, e o E1 do
          comparativo, com ela, não são intercambiáveis.
\end{itemize}

\subsection{O que este aparato transfere}

As limitações acima são específicas deste projeto. O que se propõe como
transferível é o conjunto de princípios que as camadas materializam, e que pode
ser aplicado a qualquer trabalho experimental de aprendizado de máquina aplicado
a controle:

\begin{itemize}
    \item \textbf{Toda decisão numérica tem identificador, e todo identificador
          tem teste.} O identificador é o que permite citar a decisão num
          documento, num teste e num relatório sem ambiguidade.
    \item \textbf{Toda regra registra a falha que a originou.} O custo de manter
          uma regra é sempre visível; o custo de removê-la só é visível se a
          medição que a motivou estiver escrita ao lado dela.
    \item \textbf{O critério é congelado antes da medição, e a proveniência é
          parte do veredito.} Um número correto obtido de forma não defensável
          não é reportável, e o lugar de dizer isso é o veredito, não uma nota de
          rodapé.
    \item \textbf{As ações de correção são pré-declaradas por critério.} É o que
          separa um experimento de uma busca livre, e é o que torna uma tabela de
          quatro linhas legível como demonstração.
    \item \textbf{As limitações são declaradas com o custo quantificado.} Uma
          ressalva sem número é uma ressalva que o leitor não consegue avaliar.
\end{itemize}

Nenhum desses princípios é caro em si. O que custa é a disciplina de aplicá-los
antes de a falha aparecer --- e a evidência oferecida por este capítulo é que,
neste projeto, cada um deles foi escrito \textit{depois} de uma falha medida
custar entre uma tentativa de uma hora e oito execuções reprovadas em silêncio.
